%% title: English expletive agreement
%% category: phenomena

When specTP is filled with an expletive, English verbs agree with the copular or unaccusative subject DP below them:

% There seem to be two--

\pex \dref{eng-expl-raven}
  \a There was a raven.
  \a \begin{forest}
[TP
  [DP \\ \textit{there}, roof]
  [T\1
    [T]
    [\vP
      [\lv \\ \textit{was}]
      [DP\feat{–pl} \\ \textit{a raven}, roof]]]]
\end{forest}
\xe

\pex \dref{eng-expl-brothers}
  \a There arrive the brothers.
  \a \begin{forest}
[TP
  [DP \\ \textit{there}, roof]
  [T\1
    [T]
    [\vP
      [DP\feat{+pl} \\ \textit{the brothers}, roof]
      [\lv\1
        [\lv]
        [\rt{ARRIVE}]]]]]
\end{forest}
\xe

These derivations are even simpler. At some point the relevant DP is merged with another node, and because only the DP has a number feature, that feature projects further and further. This proves that the algorithm must be \textbf{value-seeking} for number -- the rules in \Next all apply at some point in \gref{eng-expl-raven} and \gref{eng-expl-brothers}:

\pex \dref{eng-expl-rules}
  \a $\tau(\text{[–pl]}, \text{[ ]}) = \text{[–pl]}$
  \a $\tau(\text{[ ]}, \text{[–pl]}) = \text{[–pl]}$
  \a $\tau(\text{[+pl]}, \text{[ ]}) = \text{[+pl]}$
  \a $\tau(\text{[ ]}, \text{[+pl]}) = \text{[+pl]}$
\xe

There is some optionality with expletive agreement -- namely, that contracted \textit{there's} can occur with plural DPs, in certain registers. I assume this is just a separate lexical item.
